\chapter{Burning Mouth Syndrome}

\section*{Definition}
A chronic or recurrent burning sensation in the oral cavity without clinically identifiable mucosal abnormalities, commonly affecting the tongue, lips, palate, or entire mouth.

\section*{When to See a Doctor}

\begin{itemize}
 \item Burning mouth lasting more than 2 weeks with weight loss, difficulty eating, mucosal lesions, or associated systemic symptoms
\end{itemize}

\section*{How It Develops}
Primary BMS is a chronic oral pain condition diagnosed after local and systemic causes have been assessed. It may involve altered pain and taste processing, but its mechanisms are not fully understood. Dry mouth, hormonal changes, anxiety, and sleep problems may coexist or worsen symptoms; visible lesions suggest another or additional diagnosis.

\section*{Causes}
\begin{itemize}
 \item Perimenopausal or menopausal hormonal changes
 \item Xerostomia (dry mouth from medications, autoimmune, or radiation)
 \item Nutritional deficiencies (iron, zinc, B vitamins, folate)
 \item Oral candidiasis (thrush)
 \item Gastro-oesophageal reflux disease (acid reaching the mouth)
 \item Medication side effects (ACE inhibitors, antiretrovirals, chemotherapy)
 \item Anxiety, depression, or sleep disturbance, which may coexist with or amplify symptoms
 \item Allergic contact stomatitis (dental materials, flavourings)
 \item Diabetes mellitus (peripheral neuropathy)
 \item Geographic tongue or oral lichen planus
\end{itemize}

\section*{Related Symptoms}
\begin{itemize}
 \item Dry mouth
 \item Altered taste
 \item Oral thrush
 \item Mouth ulcers
 \item Bad breath
\end{itemize}


\section*{Medical Evaluation}
\begin{itemize}
 \item Your doctor will examine the entire oral mucosa, tongue, palate, and lips for erythema, atrophy, ulcers, or white plaques suggesting candidiasis or lichen planus.
 \item Laboratory evaluation includes complete blood count, iron studies, ferritin, zinc, vitamin B12, folate, and fasting glucose to screen for nutritional and metabolic causes.
 \item Swabs or other tests may be used when examination suggests infection; a normal-appearing mouth does not by itself establish candidiasis.
 \item Referral to an oral medicine specialist, dentist, or neurologist is arranged when initial workup is unrevealing and symptoms persist.
\end{itemize}

\section*{Treatment Options}
\begin{itemize}
 \item Antifungal therapy (clotrimazole troches, fluconazole) is prescribed for confirmed oral candidiasis.
 \item Nutritional supplementation (iron, zinc, vitamin B complex) corrects identified deficiencies and may improve burning symptoms over several weeks.
 \item Topical clonazepam, capsaicin, or lidocaine may be considered by an oral-medicine clinician, but evidence is limited and risks such as sedation or dependence must be discussed.
 \item Psychological support and selected neuropathic-pain medicines may help some people; treatment should be individualized rather than based solely on anxiety or depression.
\end{itemize}

\section*{Self-Care and Home Management}
\begin{itemize}
 \item Drink water frequently throughout the day and suck on sugar-free ice chips or lozenges to maintain oral moisture.
 \item Avoid irritants including hot/spicy foods, acidic foods (citrus, tomatoes), alcohol, tobacco, and mouthwashes containing alcohol.
 \item Use a soft-bristled toothbrush and a toothpaste formulated for sensitive mouths, avoiding whitening or tartar-control products.
 \item Keep a symptom diary noting foods, activities, or times of day that trigger or relieve the burning sensation.
\end{itemize}

\section*{References}
\begin{thebibliography}{99}
\bibitem{burning_mouth_syndrome:ref1} Scala A, Checchi L, Montevecchi M, et al. ``Update on burning mouth syndrome: overview and patient management.'' Crit Rev Oral Biol Med. 2003;14(4):275-291.
\bibitem{burning_mouth_syndrome:ref2} Klasser GD, Fischer DJ, Epstein JB. ``Burning mouth syndrome: a clinical update.'' J Am Dent Assoc. 2008;139(7):886-894.
\bibitem{burning_mouth_syndrome:ref3} Coculescu EC, Tovaru S, Coculescu BI. ``Burning mouth syndrome: a review on diagnosis and treatment.'' J Med Life. 2014;7(4):512-515.
\bibitem{burning_mouth_syndrome:ref4} Nasri-Heir C, Khan J, Ananthan S, et al. ``Burning mouth syndrome: a review of etiology and management.'' J Am Dent Assoc. 2021;152(5):359-369.
\bibitem{burning_mouth_syndrome:ref5} Grushka M, Epstein JB, Gorsky M. ``Burning mouth syndrome.'' Am Fam Physician. 2002;65(4):615-620.
\bibitem{burning_mouth_syndrome:ref6} Zakrzewska JM. ``Burning mouth syndrome.'' UpToDate. Wolters Kluwer; 2025.
\end{thebibliography}
